% !TEX program = lualatex
\documentclass[9pt,a4paper]{ltjsarticle}
\usepackage{amsmath,amssymb}
\usepackage{geometry}
\geometry{top=1.2cm,bottom=1.0cm,left=1.3cm,right=1.3cm}
\usepackage{hyperref}
\hypersetup{colorlinks=true,linkcolor=blue,citecolor=blue,urlcolor=blue}
\usepackage{booktabs}
\usepackage{multicol}
\usepackage{enumitem}
\usepackage{luatexja-fontspec}
\usepackage{tcolorbox}
\tcbuselibrary{skins}

\setlist{nosep,leftmargin=1.0em,topsep=0pt,partopsep=0pt}
\pagestyle{empty}
\setlength{\parskip}{0pt}
\setlength{\parindent}{0pt}
\setlength{\columnsep}{12pt}

\begin{document}

\begin{center}
{\Large\bfseries BH熱力学プログラム：6篇の研究プログラム概要}\\[2pt]
{\small 木原 範昭（Noriaki Kihara）\enspace WF System Co., Ltd.\enspace 大阪大学 基礎工学部 卒業\enspace ORCID: \href{https://orcid.org/0009-0004-6753-4020}{0009-0004-6753-4020}}
\end{center}

\vspace{-1pt}
\hrule height 0.7pt
\vspace{4pt}

\begin{tcolorbox}[colback=gray!7,colframe=gray!55,boxrule=0.4pt,arc=2pt,left=4pt,right=4pt,top=2pt,bottom=2pt]
{\small
\textbf{主張 1：4次元球への単位立方体充填の漸近係数 $c = 8/(3\pi)$}\enspace
半径 $R = 2k+1$ の4次元球内に整数中心の単位立方体を完全包含で詰めた個数 $N(k)$ について、体積不足 $\Delta(R) = V_4(R) - N(k)$ の漸近係数 $c$ を包除原理から閉形式で導出。$c = 8/(3\pi) \approx 0.84883$、数値検証で 0.024\% 一致。{\scriptsize（論文3）}\\[2pt]
\textbf{主張 2：4+1次元 Tangherlini ブラックホール熱力学との完全一致}\enspace
中心投影由来計量の自己整合的決定から $R \propto M^{1/2}$、Hawking 温度 $T_H \propto M^{-1/2}$、寿命 $\tau \propto M^{3/2}$、蒸発レート $dM/dt \propto -M^{-1/2}$ を導出。Bekenstein--Hawking エントロピーとの比 $\Delta(R)/S_{BH} \to 32/(3\pi) \approx 3.40$。{\scriptsize（論文4・5）}\\[2pt]
\textbf{主張 3：ホログラフィー非仮定の独立な経路}\enspace
't~Hooft--Susskind ホログラフィー原理、Maldacena AdS/CFT、特定の量子重力プログラム（LQG/CDT/弦理論）を\textbf{仮定しない}。中心投影と離散性という最小限の幾何学的・組合せ論的仮定のみから上記結果を導出する独立経路。自己引用ゼロ。{\scriptsize（論文1〜6）}}
\end{tcolorbox}

\vspace{2pt}

\begin{multicols}{2}
\small

\textbf{論文1：俯瞰・仮説提示}\enspace
仮定 A（BH エントロピーは幾何学的）・B（蒸発は地平面局在）・C（時空は Planck スケールで離散）から、仮説1（射影仮説）と仮説2（離散ドリフト仮説）を抽出。

\smallskip
\textbf{論文2：中心投影の数学的基盤}\enspace
$S^4(R) \subset \mathbb{R}^5$ から4次元接超平面へのグノモン投影。誘導計量 $g_{\mu\nu} = (R^2/\ell^2)(\delta_{\mu\nu} - x_\mu x_\nu/\ell^2)$ は $G_{\mu\nu} + (3/R^2)g_{\mu\nu} = 0$ を満たし、Lorentz 連続化で de Sitter、質量導入で Schwarzschild--de Sitter 形。

\smallskip
\textbf{論文3：離散構造の数学}\enspace
充填条件 $\sum (|c_i|+1/2)^2 \le R^2$、$N(k)$ を $k = 0..60$ で正確計算（$N(60) = 1{,}028{,}515{,}513$）。$\Delta(R) = (16\pi/3)R^3 - 6\pi R^2 + O(R)$。Lagrange--Jacobi 四平方表現数 $r_4(N) = 8\sigma(N)$ と接続。

\smallskip
\textbf{論文4・5：物理接続と動力学}\enspace
$\Delta(R)$ と4+1 Bekenstein--Hawking エントロピーの定数比 $32/(3\pi)$。Friedmann 類似 $E_{\mathrm{total}} \propto R^2$ から Tangherlini スケーリング全項目を回復。

\smallskip
\textbf{論文6：統合と未解決問題}\enspace
4+1次元での結果が観測 3+1 BH 物理学にどう削減されるかを主要未解決問題として精密に表明。3候補機構（KK・Randall--Sundrum/ADD・dS-CFT）。

\vspace{6pt}
\textbf{6篇のDOI一覧}

\vspace{1pt}
{\footnotesize
\begin{tabular}{@{}ll@{}}
論文1（俯瞰）  & \href{https://doi.org/10.5281/zenodo.19837588}{10.5281/zenodo.19837588} \\
論文2（射影） & \href{https://doi.org/10.5281/zenodo.19837590}{10.5281/zenodo.19837590} \\
論文3（充填） & \href{https://doi.org/10.5281/zenodo.19837592}{10.5281/zenodo.19837592} \\
論文4（エントロピー） & \href{https://doi.org/10.5281/zenodo.19837594}{10.5281/zenodo.19837594} \\
論文5（動力学） & \href{https://doi.org/10.5281/zenodo.19837596}{10.5281/zenodo.19837596} \\
論文6（統合） & \href{https://doi.org/10.5281/zenodo.19837598}{10.5281/zenodo.19837598} \\
\end{tabular}
}

\columnbreak

\textbf{4+1 Tangherlini との対応}

\vspace{1pt}
\begin{center}
\footnotesize
\begin{tabular}{@{}lll@{}}
\toprule
量 & 中心投影構成 & 4+1 Tangherlini \\
\midrule
質量--半径 & $R \propto M^{1/2}$ & $r_h \propto M^{1/2}$ \\
Hawking 温度 & $T_H \propto M^{-1/2}$ & 同上 \\
寿命 & $\tau \propto M^{3/2}$ & 同上 \\
エントロピー & $\Delta(R) \propto R^3$ & $S_{BH} \propto R^3$ \\
比 & \multicolumn{2}{l}{$\Delta/S_{BH} \to 32/(3\pi)$} \\
\bottomrule
\end{tabular}
\end{center}

\vspace{6pt}
\textbf{本研究が主張しないこと}

\vspace{1pt}
\begin{itemize}[topsep=1pt]
\item 物理的時空が 4+1 次元であること
\item 高次元時空の物理的実在
\item Bekenstein--Hawking エントロピーの第一原理導出
\item ホログラフィック・量子重力プログラムの代替
\end{itemize}

\vspace{1pt}
{\footnotesize 蒸発レート指数の「不一致」は 3+1 と 4+1 規約の混同であり、4+1 次元では完全一致。真の未解決問題は 4+1 $\to$ 3+1 削減機構の同定。}

\end{multicols}

\vspace{0pt}
\hrule height 0.4pt
\vspace{3pt}

\begin{center}
\begin{minipage}[c]{14cm}
{\footnotesize\textbf{Zenn 記事}\enspace \url{https://zenn.dev/noriaki_kihara/articles/bh-thermodynamics-projection}}\\[1pt]
{\footnotesize\textbf{note 記事（日本語）}\enspace \url{https://note.com/kiharanoriaki/n/n4524006ef175}}\\[1pt]
{\footnotesize\textbf{note 記事（英語）}\enspace \url{https://note.com/kiharanoriaki/n/n916eac778c1c}}\\[1pt]
{\footnotesize\textbf{ソースコード（GitHub）}\enspace \url{https://github.com/WurabeSeiji/ai-chat-logs-open/tree/main/BH熱力学プログラム}}\\[1pt]
{\scriptsize 全論文のPDF・LaTeX・MarkdownはZenodoに公開済み（日英対訳、CC BY 4.0）}
\end{minipage}
\end{center}

\end{document}
